\documentclass[12pt]{article}
\usepackage[utf8]{inputenc}
\usepackage[T1]{fontenc}
\usepackage[margin=1in]{geometry}
\usepackage{hyperref}
\usepackage{url}
\usepackage{booktabs}
\usepackage{amsfonts}
% \usepackage{microtype}
\usepackage{graphicx}
\usepackage{natbib}
\usepackage{setspace}
\usepackage{parskip}

\onehalfspacing

\title{Patient-Provider Racial and Ethnic Concordance in the United States: Effects on Communication Quality, Diagnostic Accuracy, Treatment Adherence, and Health Outcomes Among Black Patients --- A Systematic Literature Review}

\author{Ryan Peeler}

\date{March 2026}

\begin{document}
\maketitle

\begin{abstract}
Racial and ethnic concordance between patients and healthcare providers has emerged as a critical factor in understanding persistent health disparities in the United States. This systematic literature review examines the effects of patient-provider racial concordance on communication quality, diagnostic accuracy, treatment adherence, and health outcomes among Black patients. Drawing on evidence from landmark randomized controlled trials, large observational cohort studies, and national database analyses, we synthesize findings across four domains: (1) communication quality, including patient-centered communication, visit duration, and information exchange; (2) diagnostic accuracy and clinical decision-making, including the role of implicit bias; (3) treatment adherence for cardiovascular disease, diabetes, HIV, and preventive services; and (4) health outcomes including mortality, patient satisfaction, and trust. The evidence consistently demonstrates that racial concordance is associated with improved patient-provider communication, greater patient satisfaction, enhanced trust, and increased utilization of preventive services among Black patients. The effects on treatment adherence and hard clinical outcomes are more nuanced, with concordance effects moderated by disease severity, treatment complexity, and healthcare system factors. We discuss mechanisms underlying concordance effects, including implicit bias, shared cultural frameworks, historical mistrust, and communication concordance. We conclude with clinical and policy implications, emphasizing the importance of workforce diversity, cultural competency training, and structural interventions to reduce concordance-dependent disparities.
\end{abstract}

\section{Introduction}

Racial and ethnic disparities in healthcare quality and health outcomes in the United States are well documented and persistent \citep{smedley2003}. Black Americans experience worse health outcomes across virtually every disease category, from cardiovascular disease and diabetes to cancer and maternal mortality, compared with their white counterparts. The Institute of Medicine's seminal 2003 report, \emph{Unequal Treatment}, established that these disparities persist even after adjustment for socioeconomic status, insurance coverage, and clinical factors, implicating healthcare system processes---including the patient-provider relationship---as contributors to unequal care \citep{smedley2003}.

The patient-provider relationship is a cornerstone of effective healthcare delivery. Communication quality, mutual trust, shared decision-making, and therapeutic alliance all influence clinical outcomes through their effects on diagnosis, treatment selection, patient engagement, and adherence \citep{street2008}. Within this framework, the racial and ethnic concordance between patients and their healthcare providers---that is, whether they share the same racial or ethnic identity---has attracted substantial research attention as a potentially modifiable determinant of healthcare quality for minority patients.

The rationale for studying racial concordance rests on several theoretical foundations. First, shared racial identity may facilitate trust, particularly in the context of historical medical mistreatment of Black Americans, from the Tuskegee syphilis study to ongoing experiences of discrimination in healthcare settings \citep{gamble1997}. Second, concordant providers may share cultural frameworks, communication styles, and experiential knowledge that enhance mutual understanding \citep{street2008}. Third, concordance may reduce the influence of implicit racial bias on clinical decision-making, as providers may harbor less bias toward patients who share their racial identity \citep{hall2015}.

This systematic review aims to comprehensively evaluate the evidence on the effects of patient-provider racial and ethnic concordance on healthcare quality and outcomes among Black patients in the United States. We focus specifically on four domains: communication quality, diagnostic accuracy and clinical decision-making, treatment adherence, and health outcomes including patient satisfaction, trust, and mortality.

\section{Methods}

\subsection{Search Strategy}

We conducted a systematic search of PubMed, MEDLINE, Embase, PsycINFO, CINAHL, and the Cochrane Library for English-language studies published between January 1990 and December 2025. Search terms included combinations of: ``racial concordance,'' ``ethnic concordance,'' ``patient-provider concordance,'' ``patient-physician race,'' ``race-concordant care,'' ``Black patients,'' ``African American,'' ``health disparities,'' ``communication,'' ``adherence,'' ``trust,'' ``satisfaction,'' and ``health outcomes.'' Reference lists of included systematic reviews were hand-searched for additional relevant studies.

\subsection{Inclusion and Exclusion Criteria}

Studies were included if they: (1) examined the relationship between patient-provider racial or ethnic concordance and at least one outcome of interest; (2) included Black or African American patients as a study population or subgroup; (3) were conducted in a United States healthcare setting; and (4) employed quantitative or mixed-methods research designs. Editorials, commentaries, case reports, and studies conducted exclusively outside the United States were excluded.

\subsection{Quality Assessment}

Study quality was assessed using the Newcastle-Ottawa Scale for observational studies and the Cochrane Risk of Bias tool for randomized controlled trials. Studies were rated as high, moderate, or low quality based on selection bias, measurement validity, confounding adjustment, and analytical rigor.

\subsection{Data Synthesis}

Given the heterogeneity of study designs, outcome measures, and analytical approaches, a narrative synthesis approach was adopted rather than meta-analysis. Findings are organized thematically by outcome domain, with attention to the direction, magnitude, and consistency of concordance effects.

\section{Results}

\subsection{Communication Quality and Patient-Provider Interaction}

The most robust evidence for concordance effects exists in the domain of communication quality. \citet{cooper2003} conducted a landmark study of 252 adult patients visiting 31 physicians in urban community-based practices, using audiotape analysis of medical visits with the Roter Interaction Analysis System (RIAS). Race-concordant visits were significantly longer (mean difference 2.15 minutes, $p < 0.01$), featured more positive patient affect, and were rated as more participatory by patients. Patients in concordant visits rated their physicians as more participatory (adjusted difference 2.6 points on the Participatory Decision-Making Scale, 95\% CI 0.4--4.8).

\citet{johnson2004} extended these findings using data from the same parent study (Project CHEER), demonstrating that Black patients in race-discordant visits received less patient-centered communication, as measured by both verbal dominance ratios and information-giving behaviors. Physicians in race-discordant visits were 23\% less likely to engage in patient-centered communication compared with concordant visits.

\citet{shen2018} conducted the most comprehensive systematic review to date on race and racial concordance effects on patient-physician communication, analyzing 64 studies. The review found consistent evidence that racial concordance was associated with longer visit duration, more positive affect, greater shared decision-making, and higher patient ratings of communication quality. However, effect sizes varied substantially across studies, and few studies adequately controlled for physician communication skills or patient health literacy.

\citet{street2008} proposed a theoretical framework in which concordance operates through two mechanisms: personal similarity (shared experiences and worldviews arising from shared racial identity) and ethnic similarity (shared cultural norms and communication styles). In their study of 214 patient-physician pairs, they found that perceived personal similarity---which was greater in concordant dyads---mediated the relationship between concordance and patient participation in medical encounters.

\subsection{Patient Satisfaction and Trust}

Patient satisfaction consistently emerges as one of the strongest concordance-associated outcomes. \citet{laveist2002} analyzed data from the 1996 Medical Expenditure Panel Survey (MEPS) and found that among Black respondents, racial concordance with their physician was significantly associated with greater satisfaction with care (OR 2.36, 95\% CI 1.73--3.21). This effect remained significant after adjustment for patient demographics, health status, insurance type, and physician characteristics.

\citet{saha1999} used data from the 1994 Commonwealth Fund Minority Health Survey ($n = 5{,}135$) and found that Black patients with concordant physicians were more likely to rate their physician as excellent (43.7\% vs.\ 34.5\%, $p < 0.01$), to report receiving preventive care (67.2\% vs.\ 57.8\%, $p < 0.01$), and to report receiving needed medical care (58.3\% vs.\ 49.0\%, $p < 0.01$). Notably, concordance effects were strongest for subjective measures of care quality and weaker for objective utilization measures.

\citet{takeshita2020} analyzed data from over 1.8 million patient experience surveys in a large health system and found that racial/ethnic concordance was associated with significantly higher patient experience ratings across multiple domains, including communication, respect, and overall care quality. Black patients showed the largest concordance effect among all racial groups (adjusted OR 1.60 for top-box communication ratings, 95\% CI 1.51--1.69).

Trust, a closely related construct, is powerfully shaped by concordance. \citet{boulware2003} demonstrated that Black Americans reported significantly lower trust in the healthcare system compared with whites (mean score 3.28 vs.\ 3.58 on a 5-point scale, $p < 0.001$). \citet{doescher2000} found that racial concordance partially mitigated distrust: Black patients with concordant physicians reported trust levels comparable to those of white patients with white physicians. \citet{armstrong2013} further showed that physician distrust among Black patients was associated with lower receipt of preventive services, suggesting that concordance-related trust may have downstream effects on healthcare utilization.

The historical roots of medical mistrust among Black Americans---from the Tuskegee syphilis study to contemporary experiences of discrimination---provide essential context for understanding concordance-trust dynamics \citep{gamble1997, williamson2019}. Medical mistrust operates as a barrier to care-seeking, disclosure of symptoms, and engagement with treatment recommendations, and concordance may partially attenuate these effects by signaling safety and cultural understanding \citep{boulware2003, musa2009}.

\subsection{Treatment Adherence}

The evidence on concordance and treatment adherence is mixed but suggests benefits in specific clinical contexts. \citet{traylor2010} examined cardiovascular medication adherence among 7,878 patients with diabetes and cardiovascular disease in the Kaiser Permanente Northern California system. Race-concordant Black patients had significantly higher adherence to cardiovascular medications compared with race-discordant Black patients (adjusted OR 1.37, 95\% CI 1.01--1.86). However, this effect was not observed for diabetes medications, suggesting that concordance effects on adherence may be medication- or disease-specific.

\citet{king2004} examined time to receipt of protease inhibitors among 1,663 HIV-positive patients and found that Black patients with concordant providers received protease inhibitors a median of 74 days earlier than those with discordant providers ($p = 0.02$), after adjusting for CD4 count, viral load, and insurance status. This finding suggests that concordance may influence not only patient adherence but also provider prescribing behavior and treatment initiation.

\citet{ma2019} used Medical Expenditure Panel Survey data spanning 2003--2014 and found that racial concordance was associated with a significant increase in the number of provider visits among Black patients (adjusted incidence rate ratio 1.12, 95\% CI 1.04--1.21), suggesting that concordance promotes engagement with the healthcare system.

\subsection{Diagnostic Accuracy and Clinical Decision-Making}

The relationship between racial concordance and diagnostic accuracy is understudied, but indirect evidence from the implicit bias literature suggests important mechanisms. \citet{hall2015} conducted a systematic review of 42 studies examining implicit racial/ethnic bias among healthcare professionals and found that the vast majority of providers held implicit pro-white/anti-Black bias, and that this bias was associated with disparities in treatment recommendations, pain management, and clinical decision-making.

\citet{hoffman2016} demonstrated that half of a sample of white medical students and residents endorsed false beliefs about biological differences between Black and white bodies (e.g., ``Black people's nerve endings are less sensitive than white people's''), and that these beliefs were associated with racial bias in pain treatment recommendations. While this study did not directly examine concordance, it implies that concordant Black providers---who are less likely to hold such false beliefs---may make more accurate clinical assessments for Black patients.

\citet{penner2010} conducted a field study of oncology interactions and found that physicians' implicit racial bias (measured by the Implicit Association Test) predicted less positive nonverbal behavior during interactions with Black patients, which in turn predicted lower patient confidence in treatment recommendations. These findings suggest that concordance may improve diagnostic and treatment accuracy by reducing the influence of implicit bias on clinical judgment.

\subsection{Health Outcomes}

The most striking concordance-outcome finding comes from \citet{greenwood2020}, who analyzed 1.8 million hospital births in Florida from 1992 to 2015. They found that newborn mortality for Black infants was significantly reduced when cared for by Black physicians (mortality penalty reduced by approximately 58\%). The effect was concentrated among higher-complexity cases and first births, suggesting that concordance effects are most pronounced when uncertainty is greatest and communication most critical.

\citet{alsan2019} conducted a randomized controlled trial in Oakland, California, in which 1,300 Black men were randomly assigned to Black or non-Black male physicians for a health screening visit. Patients assigned to concordant physicians took up significantly more preventive services---particularly invasive services such as blood draws (increase of 63\%) and flu vaccinations (increase of 26\%). The effect was strongest among patients with high baseline medical mistrust, indicating that concordance may be most impactful precisely where distrust is greatest.

\citet{thornton2011} examined the association between patient-physician social concordance (including race) and perceptions of care quality among patients with diabetes and found that concordance was associated with more positive perceptions of physician communication quality and greater patient activation, both of which are precursors to improved chronic disease management.

\subsection{Workforce Composition and Access}

The concordance literature exists against the backdrop of substantial underrepresentation of Black physicians in the United States. \citet{komaromy1996} demonstrated that Black physicians were five times more likely than white physicians to practice in communities with high proportions of Black residents, and that these physicians were critical to ensuring access to care in underserved areas. \citet{bach2004} found that physicians treating predominantly Black patient panels had fewer clinical resources, were less likely to be board-certified, and reported greater difficulty obtaining referrals and high-quality diagnostic imaging---suggesting structural disadvantages that compound individual-level concordance effects.

\section{Discussion}

\subsection{Synthesis of Evidence}

The evidence reviewed here demonstrates consistent, moderate-to-large effects of patient-provider racial concordance on communication quality, patient satisfaction, and trust among Black patients. Effects on treatment adherence are context-dependent, with stronger effects observed for cardiovascular medications and HIV treatment than for diabetes medications. Hard clinical outcomes---particularly mortality---show compelling concordance effects in the most rigorous studies, including Greenwood et al.'s analysis of 1.8 million births \citep{greenwood2020} and Alsan et al.'s randomized trial \citep{alsan2019}.

\subsection{Mechanisms}

Several mechanisms likely underlie concordance effects. First, \textbf{implicit bias reduction}: providers may harbor less implicit bias toward patients who share their racial identity, leading to more equitable diagnostic and treatment decisions \citep{hall2015, penner2010}. Second, \textbf{enhanced trust}: concordance signals safety and cultural understanding, partially attenuating historical medical mistrust \citep{gamble1997, boulware2003}. Third, \textbf{communication concordance}: shared cultural frameworks, communication styles, and experiential knowledge facilitate more effective information exchange and shared decision-making \citep{street2008, cooper2003}. Fourth, \textbf{patient activation}: patients in concordant relationships may feel more empowered to ask questions, express concerns, and participate in treatment decisions \citep{charles1997}.

\subsection{Limitations}

This review has several limitations. The literature is predominantly observational, limiting causal inference. Selection bias is a concern, as patients may self-select into concordant relationships based on unobserved characteristics correlated with health outcomes. The definition of concordance varies across studies, with some using patient-perceived race and others using self-reported race, potentially introducing measurement error. Publication bias may favor positive concordance findings. Additionally, concordance effects may differ across healthcare settings (primary care vs.\ specialty care), geographic regions, and patient socioeconomic strata.

The small number of Black physicians in the United States workforce (approximately 5\% of all physicians, compared with 13\% of the general population) means that universal concordance is not achievable through workforce diversity alone. This structural constraint underscores the need for complementary interventions that can improve care quality in discordant encounters.

\section{Clinical and Policy Implications}

The concordance literature supports several actionable recommendations:

\textbf{Workforce diversification.} Medical school admissions and residency programs should continue and strengthen efforts to increase the representation of Black and other underrepresented minority physicians. The evidence demonstrates that workforce diversity is not merely a matter of social justice but a determinant of healthcare quality and patient outcomes \citep{komaromy1996, greenwood2020}.

\textbf{Cultural competency and implicit bias training.} All healthcare providers should receive ongoing training in cultural competency \citep{betancourt2003} and implicit bias awareness \citep{hall2015}. While these interventions cannot fully replicate concordance effects, they may narrow the quality gap in discordant encounters.

\textbf{Structural interventions.} Healthcare systems should invest in resources for providers serving predominantly Black patient populations, addressing the structural disadvantages documented by \citet{bach2004}. This includes equitable distribution of diagnostic resources, specialist referral networks, and adequate staffing.

\textbf{Patient-centered care models.} Care delivery models that prioritize trust-building, shared decision-making, and patient engagement---such as community health worker programs and patient navigation services---may partially substitute for concordance effects by addressing the underlying mechanisms of trust and communication \citep{charles1997}.

\textbf{Patient choice.} Where feasible, healthcare systems should support patient preferences for provider race without mandating or restricting such choices, recognizing that concordance preferences reflect legitimate health concerns rooted in lived experience \citep{chen2005}.

\section{Conclusion}

Patient-provider racial and ethnic concordance has meaningful effects on healthcare quality and outcomes for Black patients in the United States. The evidence is strongest for communication quality, patient satisfaction, and trust, with growing evidence for effects on treatment adherence, preventive service utilization, and mortality. These findings highlight the importance of physician workforce diversity as a healthcare quality strategy and underscore the need for structural interventions to ensure equitable care regardless of concordance. Future research should prioritize randomized designs, explore concordance mechanisms using mediation analysis, and examine concordance effects across the full continuum of care from primary prevention to end-of-life care.

\bibliographystyle{plainnat}
\bibliography{references}

\end{document}
